\section{Transformer attention-only de una capa: primero se expresa el modelo como rutas aditivas}

Para obtener un álgebra exacta, Olah retira temporalmente las MLP, la LayerNorm y los bias, y conserva solo una capa de multi-head attention, la residual connection, el token embedding y el unembedding. No es la arquitectura completa de un LLM moderno, sino un banco de pruebas: si ni siquiera este sistema puede descomponerse, será todavía más difícil explicar los modelos grandes.

\lecturefigure{slide-11-one-layer-simplifications.jpg}{Supuestos del análisis: one layer, attention-only, sin LayerNorm y sin bias.}{Video oficial de Stanford Online 00:15:32--00:15:59.}

\begin{warningbox}{Función y límites del toy model}
El modelo simplificado permite escribir cada path como un producto de matrices y llevar una accounting casi completa. Puede revelar reusable concepts, pero no demuestra que las mismas fórmulas sigan dominando el comportamiento de la misma manera en modelos con MLP, RMSNorm, RoPE, GQA y cientos de capas.
\end{warningbox}

\subsection{Forma del residual stream}

Sea la token one-hot matrix $T\in\mathbb{R}^{n\times |V|}$, el embedding $W_E\in\mathbb{R}^{|V|\times d}$ y el unembedding $W_U\in\mathbb{R}^{d\times |V|}$. El residual stream inicial es
\[
X_0=T W_E.
\]
El attention-only model de una capa se escribe como
\[
X_1=X_0+\sum_{h\in H}A^h(X_0)X_0W_{OV}^h,
\qquad
\mathrm{logits}=X_1W_U.
\]

\begin{itemize}
  \item $A^h\in\mathbb{R}^{n\times n}$: el attention pattern de la head $h$;
  \item $W_{OV}^h=W_V^hW_O^h\in\mathbb{R}^{d\times d}$: matriz compuesta de value/output;
  \item la residual identity path conserva la representación original de la posición actual;
  \item la salida de cada head se escribe de vuelta, por suma, en el shared residual stream.
\end{itemize}

\begin{knowledgebox}{Términos clave: niveles de objetos en el análisis de circuitos}
\small
\begin{tabularx}{\textwidth}{L{0.22\textwidth}L{0.28\textwidth}X}
\toprule
Término & Pregunta que responde & Función en esta clase \\
\midrule
Residual stream & ¿Dónde reside la representación actual que comparten todas las capas? & Funciona como un espacio de trabajo común; la representación de entrada, cada cabeza de atención y las capas posteriores leen o escriben, mediante sumas, el mismo conjunto de vectores de posición. \\
Attention pattern $A$ & ¿De qué source position a qué destination position fluye la información? & Es un position-to-position routing; no indica directamente el contenido de la feature que se transporta. \\
QK circuit & ¿Por qué la posición actual atiende a cierta posición anterior? & $W_{QK}=W_QW_K^{\top}$ define la query--key compatibility; es la regla parametrizada de ``where to read''. \\
OV circuit & Una vez leída la información, ¿qué escribe la head? & $W_{OV}=W_VW_O$ convierte la source feature en una residual update y puede plegarse además hacia los output-token logits. \\
Path composition & ¿Por qué componentes pasó un resultado? & Despliega la acción sucesiva del residual, de las heads de la primera capa y de las de la segunda, y distingue direct, single-head y composed paths. \\
Virtual head & ¿La combinación de dos heads reales forma una función nueva? & Es un effective circuit que surge del producto de matrices entre capas, no un conjunto de parámetros nuevo dentro del modelo. \\
Induction head & ¿Cómo implementa el modelo $[A][B]\ldots[A]\mapsto[B]$? & La capa anterior mueve previous-token information y la capa posterior copia la continuation después de un contexto histórico similar. \\
Ablation & ¿El circuit candidato solo está correlacionado o contribuye causalmente al comportamiento? & Se retira o se reemplaza cierta head/path y se mide el cambio del ICL score; la conclusión solo abarca la intervención y la metric utilizadas. \\
\bottomrule
\end{tabularx}
\end{knowledgebox}

\lecturefigure{slide-12-attention-only-model-algebra.jpg}{Álgebra del embedding, las heads, el residual y el unembedding en el Transformer attention-only.}{Video oficial de Stanford Online 00:16:00--00:18:13.}

\teachervoice{El ponente considera que la notación tradicional $Q,K,V,O$ tiende a ocultar la estructura de la head. Una descomposición más útil es esta: $A$ decide de dónde a dónde se mueve la información, y $W_{OV}$ decide qué contenido se mueve exactamente.}

\subsection{La attention head como operador sobre dos ejes independientes}

Si se separan el sequence-position mixing y el feature-channel mapping, la head puede verse como un tensor/Kronecker product:
\[
h(X)=A^h(X)XW_{OV}^h
\quad\Longleftrightarrow\quad
\mathcal{H}^h=A^h\otimes W_{OV}^h.
\]

\begin{itemize}
  \item multiplicar por la izquierda por $A^h$: weighted sum entre positions;
  \item multiplicar por la derecha por $W_{OV}^h$: lectura y escritura en el residual feature space;
  \item los parámetros y la semántica de ambos factores pueden examinarse por separado;
  \item pero $A^h$ depende de la entrada, por lo que la head completa sigue siendo no lineal.
\end{itemize}

\lecturefigure{slide-13-attention-head-tensor-product.jpg}{Una sola attention head: position mixing $A^h$ y content map $W_{OV}^h$.}{Video oficial de Stanford Online 00:18:14--00:18:31.}

\lecturefigure{slide-14-attention-layer-sum.jpg}{Una capa de attention: identity residual path más todos los head operators.}{Video oficial de Stanford Online 00:18:32--00:19:25.}

En esta figura, $I$ no es un adorno, sino la residual path. Permite que la representación original del token llegue directamente al unembedding, y también hace que, al desplegar varias capas, surja de forma natural la path combinatorics de «elegir una head de cierta capa o saltarla».

\subsection{Fold de embedding y unembedding}

Ya se descompuso una head en position routing y feature mapping, pero al lector le interesan en última instancia los output-token logits, no las residual coordinates abstractas. Por eso, esta sección pliega el embedding de la entrada y el unembedding de la salida dentro de cada ruta, y plantea directamente la pregunta: una vez leído cierto source token, ¿qué candidate next tokens sube o baja? Esta representación token-to-token también establece una interfaz uniforme para el examen posterior de los QK/OV circuits.

Al combinar $W_E$, $W_U$ y cada head, se obtiene directamente un vocabulary-to-vocabulary map:
\[
\mathrm{logits}
=T W_EW_U
+\sum_h A^h(T)
T\underbrace{W_EW_{OV}^hW_U}_{M_{OV}^h}.
\]

\begin{itemize}
  \item $W_EW_U$: direct path; almacena principalmente bigram-like statistics;
  \item $M_{OV}^h\in\mathbb{R}^{|V|\times |V|}$: el efecto de esa head sobre los output-token logits después de ver cierto source token;
  \item $A^h$ decide desde qué source position se lee;
  \item por lo tanto, la unidad interpretable del one-layer model pasa a ser «attention pattern × token-to-token map».
\end{itemize}

\lecturefigure{slide-15-fold-embedding-unembedding.jpg}{Después de plegar el embedding/unembedding, el modelo se divide en la direct/bigram path y las head paths.}{Video oficial de Stanford Online 00:19:26--00:19:35.}

\lecturefigure{slide-16-fixed-pattern-linearity.jpg}{Observación central: si se fijan los attention patterns, el one-layer attention-only model es lineal respecto a la entrada.}{Video oficial de Stanford Online 00:19:36--00:20:07.}

\begin{importantbox}{Qué significa «lineal una vez fijado el pattern»}
El modelo completo sigue dependiendo de la entrada de forma no lineal a través de la softmax attention; pero, para un prompt concreto o para una clase pequeña de patterns estables, podemos tratar $A^h$ temporalmente como constante y analizar la content path con álgebra lineal. El Mechanistic analysis suele aprovechar este condicionamiento local, en lugar de afirmar que el Transformer es globalmente lineal.
\end{importantbox}

\teachervoice{La metodología de Olah es esta: siempre que una parte pueda fijarse y convertirse en lineal, se obtiene una herramienta de explicación poderosa. Lo esencial es registrar con honestidad «qué se fijó», y no hacer pasar la conditional linearity por global linearity.}

\subsection{Resumen de la sección}

Un Transformer attention-only de una capa puede desplegarse en la residual direct path más un término aditivo por cada head. La perspectiva del tensor product separa «dónde se lee y escribe» de «qué se lee y escribe»; al plegar el embedding/unembedding, el efecto de cada head sobre los vocabulary logits puede examinarse directamente.
